\section{Related Work}
\label{sec:related-work}

\subsection{Graph-based Fraud Detection on Financial and Blockchain Graphs}
\label{subsec:rw-graph-fraud}

Graph-based fraud detection models represent transactions, accounts, users, or actors as nodes and use their transfer, interaction, or behavioral relations as edges. Early and recent graph fraud detectors usually start from the observation that fraud signals are difficult to identify from isolated attributes, because suspicious behaviors are often reflected in local neighborhoods, multi-hop paths, heterophilic interactions, and class-imbalanced relational patterns. PC-GNN addresses the imbalance issue by combining a label-balanced sampler with a neighbor selection module, so that the aggregation process is less dominated by majority benign nodes and more sensitive to fraud-related neighbors \cite{liu2021pcgnn}. PMP focuses on the mixture of homophily and heterophily in fraud graphs. It partitions neighbors into different label-aware groups and applies node-specific aggregation functions to distinguish homophilic and heterophilic messages, which helps prevent the representations of rare fraudulent nodes from being overwhelmed by benign neighbors \cite{zhuo2024pmp}. HOGRL further targets disguised multi-hop fraud relations. It constructs high-order transaction graphs, learns pure representations at different orders, and uses a mixture-of-experts attention module to select useful high-order signals for credit card fraud detection \cite{zou2024hogrl}. Grad studies adaptive camouflage, where fraudsters mimic benign behavioral features. It combines supervised graph contrastive learning with a guided relation diffusion generator to synthesize auxiliary homophilic relations, thereby amplifying weak fraud signals during aggregation \cite{yang2025grad}.

Recent graph fraud detection studies also pay more attention to semantic cues and interpretability. CGNN studies graph fraud detection with extremely limited labels. It performs category semantic decomposition by separating node features into normal and abnormal components, and then applies context-aware aggregation and denoising attention to exploit limited supervision more effectively \cite{li2025cgnn}. SEFraud integrates prediction and explanation by learning feature masks and edge masks inside a customized heterogeneous graph transformer, and introduces a triplet loss to improve mask learning for self-explainable fraud detection \cite{li2024sefraud}. These methods show that effective fraud detection often requires more than ordinary message passing: the model needs to handle imbalance, camouflage, heterophily, high-order structure, semantic ambiguity, and explanation requirements.

However, the above methods are mainly designed to improve fraud recognition on a given graph or under a training protocol where the current graph structure can be directly used for detection. They do not explicitly study the task-only temporal setting where later tasks cannot store, access, or replay historical nodes, edges, adjacency matrices, or subgraphs. When blockchain transaction relations and actor behaviors evolve over time, a detector trained only for a static or single-window graph can fit the current snapshot, yet it lacks a mechanism for preserving fraud-discriminative semantics across temporal snapshots. TASD-CL uses an explicit normal/abnormal semantic interface and further studies how such fraud-discriminative semantics can be preserved under task-only snapshot constraints.

\subsection{Temporal Fraud Detection and Dynamic Graph Learning}
\label{subsec:rw-temporal-dynamic}

Temporal financial fraud detection emphasizes that graph construction and message passing must respect chronological order. BRIGHT is a representative real-time fraud detection framework. It points out two practical issues in production fraud detection: future transactions must not be used to predict past risks, and graph query plus GNN inference latency can be too high for online services. To address these issues, BRIGHT designs a two-stage directed graph transformation to restrict messages to historical payment transactions, and uses a Lambda Neural Network to decouple batch entity embedding computation from real-time transaction prediction \cite{lu2022bright}. ZipZap studies blockchain fraud detection from a large-scale transaction sequence perspective. It improves the efficiency of language-model-based transaction modeling through frequency-aware compression and an asymmetric training paradigm, which reduces the cost of training on large blockchain transaction data \cite{hu2024zipzap}. These works support the practical motivation of our setting: temporal causality and deployment efficiency are central constraints in financial and blockchain fraud detection.

Dynamic graph learning provides another related line of work. Many dynamic graph models maintain temporal histories, recurrent states, memory modules, or historical neighborhoods to capture graph evolution. MemMap proposes an adaptive latent memory space for dynamic graph learning, where memory cells correspond to evolving topics of the dynamic graph and are updated as the graph changes \cite{ji2024memmap}. RTGCN studies temporal graphs from the perspective of structural roles. It uses a structural role-based GNN module to capture global structural similarity and a structural role-based GRU to model role evolution along the temporal dimension \cite{du2024rtgcn}. These methods are useful for modeling evolving graphs, because they preserve temporal signals through historical states, memory, or role evolution.

The task-only blockchain fraud setting studied in this paper has a stricter access constraint. At task $t$, the model can use only the current time-window snapshot for training. Historical graph topology, multi-hop neighborhoods, adjacency matrices, and subgraphs are unavailable in later tasks. This rules out methods whose temporal modeling depends on revisiting historical structures or maintaining historical node-level graph context. TASD-CL therefore treats temporal evolution as semantic concept drift in normal behavioral semantics and abnormal fraud semantics. Its goal is to preserve fraud-discriminative semantics in a structure-free way while adapting to the current snapshot.

\subsection{Continual Graph Learning and Structure-free Semantic Preservation}
\label{subsec:rw-cgl}

Continual learning studies sequential tasks where a model must learn new data while maintaining useful knowledge acquired from previous tasks. Representative generic methods include regularization-based learning, distillation-based learning, and replay-based learning. EWC estimates parameter importance and penalizes changes to important parameters \cite{kirkpatrick2017ewc}. LwF uses the previous model as a teacher and transfers old output behavior through distillation \cite{li2017lwf}. Experience replay stores previous examples and reuses them during later training to stabilize model behavior \cite{robins1995catastrophic}. These methods provide useful generic baselines, but they were originally developed for ordinary Euclidean data or standard task sequences. Directly attaching them to GNNs does not specify which fraud semantics should be preserved, and replay-based variants can conflict with task-only graph constraints if historical graph samples or neighborhoods are stored.

Continual graph learning further considers the structural dependency of graph data. A central difficulty is that a node cannot be replayed independently from its graph context, since GNN prediction usually depends on its computation neighborhood. TEM explicitly analyzes this issue in expanding networks. It shows that directly applying memory replay to continual graph learning may cause memory explosion because representative nodes need associated topological neighborhood structures; it therefore compresses ego-subnetworks into topology-aware embeddings under a parameter-decoupled GNN framework \cite{zhang2024tem}. Continual graph learning surveys also categorize existing methods into regularization, replay, distillation, architecture, and data-centric strategies, confirming that graph structure makes continual learning substantially different from continual learning on independent samples \cite{wang2024cglsurvey}.

Our work shares the sequential learning motivation with continual graph learning, while the target problem is task-only temporal blockchain fraud detection. The model cannot store historical nodes, edges, adjacency matrices, or subgraphs, and the main challenge is semantic concept drift between evolving normal behaviors and evolving fraud strategies. TASD-CL therefore preserves structure-free fraud-discriminative semantics at three levels. SSF protects semantic decomposition and routing parameters according to their semantic roles. SPC stores compact Gaussian prototypes in normal and abnormal semantic subspaces instead of historical graph structures. SCD distills high-confidence normal and abnormal semantics from the previous model on the current snapshot. This design differs from ordinary graph replay and generic continual learning, because the preserved knowledge is explicitly tied to normal and abnormal fraud semantics under the task-only snapshot constraint.